% Proof module for the mirrored-guard presentation.
% Prerequisites and labels are in independent_work/mirror_basis.tex.
% The zero-source argument adapts the pinned source_dispatch proof,
% commit 1517d9d79ce33f631e0f82960c849427cbba23bf, with attribution.

\section{Why the mirrored guard supports source selection}
\label{compat:source-compatibility}

The source restriction is syntactic: a permitted transport source is any
term other than the single literal $0$. In particular, a variable or an
$I$-headed term whose numerical value is zero is permitted. The numerical
cases and equation induction established above were proved under precisely
this restriction. Their derivations did not substitute into an existing
proof at a variable source. This point makes the following transfer
noncircular.

\begin{theorem}[Restoration of zero-source transport]
\label{compat:zero-restored}
The mirrored-guard calculus with transport restricted to sources
$s\not\equiv0$ derives exactly the same sequents as its version with
unrestricted transport. The formula language, contexts, computations,
and introduction and elimination rules are unchanged.
\end{theorem}

We supply the zero-source construction, adapting the argument of the
pinned \texttt{source\_dispatch} package~\cite{dispatch}. Its application here is new:
the unrestricted atomic case/progress language of that package has been
replaced by the mirrored guarded language. Lemmas~\ref{mir:cases}
and~\ref{mir:equation-induction} supply the numerical rules that the
construction requires.

\subsection{A singleton-one test using safe sources}

Use the abbreviations
\[
\begin{aligned}
 C(v,a,b)&=I(v,a;u.b),& 1&=S0,&2&=S1,\\
 Z(v)&=C(v,1,0),& K(v)&=C(v,0,1),& H(v)&=C(v,1,2),\\
 N(n)&=I(n,0;u.Z(u)),&
 R(n)&=I(n,0;u.H(u)),& \ell(n)&=N(R(n)).
\end{aligned}
\]
Every displayed update variable is fresh for the fixed parameters.
In this subsection all case distinctions are performed in the empty
retained context. The results may subsequently be weakened into a context
containing a zero-source equation.

\begin{lemma}[Uniform test and filter equations]
\label{compat:zero-tests}
At every displayed target, the restricted calculus proves
\[
 H(v)=S K(v),\qquad R(1)=1,\qquad R(SSv)=2,
\]
\[
 N(1)=1,\qquad N(2)=0,\qquad \ell(Sv)=Z(v),
\]
and
\[
                 v=C(Z(v),v,0).
                 \tag{Filter}\label{compat:filter}
\]
These are uniform derived rules at arbitrary terms, without an assumption
that numerical substitution into restricted derivations is admissible.
\end{lemma}
\begin{proof}
For $H(v)=S K(v)$ use numerical cases. At zero both sides compute to
$1$; at a successor both compute to $2$. A computation may be reversed
as an equation because its original left side is $I$-headed. Joining
the resulting equations through $1$ or $2$ uses a permitted numeral
source. Thus this case argument has no zero-source transport.

Finite forward computations give $R(1)=1$, $N(1)=1$, and $N(2)=0$.
For example, $N(1)=Z(N(0))=Z(0)=1$, and the next iterate gives
$N(2)=Z(1)=0$. Every replaced subterm is the left side of a computation
or an already derived $I$-headed equation.

The first identity and the successor computation give
\[
 R(Sv)=H(R(v))=S K(R(v)).
\]
Replace the $I$-headed subterm $R(Sv)$ in
$R(SSv)=H(R(Sv))$. The result computes to $2$.
Hence $\ell(1)=1$ and $\ell(SSv)=0$ by the computed values of $N$.
Numerical cases now give $\ell(Sv)=Z(v)$: the zero case joins the two
values $1$, while the successor case joins the two $I$-headed terms
known to be zero using Lemma~\ref{mir:zero-join}.

For the filter equation use numerical cases on its displayed variable.
At zero, the right side computes forwards to zero; at a successor it
computes forwards to that successor. Reverse each resulting equation
at its $I$-headed source. This proves the required endpoints and then
the identity at any target. The argument uses no retained assumption
depending on the case variable.
\end{proof}

\subsection{Zero-source symmetry}

\begin{lemma}
\label{compat:zero-symmetry}
For every numerical term $t$, the restricted calculus derives
\[
                        0=t\ \vdash\ t=0.
\]
\end{lemma}
\begin{proof}
If $t$ is literally zero, use reflexivity. Otherwise retain $0=t$,
keep $t$ fixed, and abbreviate
\[
                       A(n)=I(n,t;u.Su).
\]
We prove the equation $n=A(n)$ by the already established equation
induction rule, Lemma~\ref{mir:equation-induction}.

At zero, $A(0)=t$ is a computation. Its safe symmetry gives $t=A(0)$.
Transport the retained $0=t$ at the permitted source $t$ to obtain
$0=A(0)$.

For the step, retain $z=A(z)$ at the same fresh parameter $z$.
Its source is the variable $z$, hence permitted. Safe successor
congruence gives $Sz=S A(z)$. Computation and safe symmetry give
$S A(z)=A(Sz)$; transport at the successor-headed source $S A(z)$
gives $Sz=A(Sz)$. In particular, this step proof is used at its
original parameter: it is never specialized to zero.

Induction at target $1$ yields $1=A(1)$. Forward computations give
$A(1)=St$, so transport at source $A(1)$ yields $1=St$.
Its source is the permitted numeral $1$. Congruence for $\ell$ yields
$\ell(1)=\ell(St)$. Replace the $I$-headed terms using
$\ell(1)=1$ and $\ell(St)=Z(t)$, obtaining $1=Z(t)$, and safely reverse
this equation to obtain $Z(t)=1$.

Weaken the unconditional filter equation~\eqref{compat:filter} into the
retained context. Transport at the $I$-headed source $Z(t)$ turns its
right side into $C(1,t,0)$. This term computes to zero, again at an
$I$-headed source. The result is $t=0$.
\end{proof}

\subsection{All equality reasoning and all formula templates}

\begin{lemma}[Unrestricted numerical equality algebra]
\label{compat:numerical-equality-restored}
The restricted calculus admits symmetry, transitivity, and numerical
congruence for all equations, without a source restriction.
\end{lemma}
\begin{proof}
Symmetry at a nonzero syntactic source was already available; at zero
use Lemma~\ref{compat:zero-symmetry}. For congruence from $0=t$ with
$t\not\equiv0$, first reverse it, apply permitted congruence at source
$t$, and use the now unrestricted symmetry. The case $t\equiv0$ is
reflexivity.

Transitivity through a middle term other than literal zero uses one
permitted transport. For $r=0$ and $0=t$, either an endpoint is literally
zero and a premise is already the conclusion, or both endpoints are
nonzero syntactic terms. In the latter case reverse $0=t$ and apply
the safe zero-join lemma to $r=0$ and $t=0$. None of these arguments
uses the unrestricted rule being derived.
\end{proof}

\begin{lemma}[Unrestricted canonical-template transport]
\label{compat:full-transport}
For every canonical formula template $F(v)$,
\[
 \frac{\Gamma\vdash s=t\qquad\Gamma\vdash F(s)}
      {\Gamma\vdash F(t)}
\]
is admissible in the source-restricted mirrored calculus.
\end{lemma}
\begin{proof}
Induct externally on the number of binders in the finite template,
uniformly in contexts and fixed parameters. For an equation use
Lemma~\ref{compat:numerical-equality-restored}: congruence on its two
terms, symmetry, and transitivity give the desired replacement.

For $F(v)=D_{x,y}B(v,x,y)$, eliminate $F(s)$ at a fresh witness $a$,
retaining its clause $Q_yB(s,a,y)$, where $Q_yC$ here abbreviates
the reusable clause $\langle C(0);y.C(y)\rangle$ rather than a new
object-formula constructor. To introduce $F(t)$, assume
$B(t,a,0)$ and seek $B(t,a,Sz)$. The induction hypothesis, using
symmetry of $s=t$, gives $B(s,a,0)$. Use the old clause at $Sz$ and
transport forwards by the induction hypothesis to obtain
$B(t,a,Sz)$.

For a compound matrix this is the primitive case introduction.
For an atomic canonical matrix it is a derived case introduction:
its zero instance is a theorem by safe symmetry of I0; cut that theorem
into the displayed case premise, ignore the local progress hypothesis,
and apply atomic progress introduction. Thus the construction is valid
for every canonical matrix. Close the introduction and the elimination.
All eigenparameters are fresh for the equality and retained data.
\end{proof}

\begin{proof}[Proof of Theorem~\ref{compat:zero-restored}]
Each restricted proof is an unrestricted proof. Conversely, transform an
unrestricted proof recursively and replace each transport node by
Lemma~\ref{compat:full-transport}. Other inference nodes are unchanged.
The contexts and conclusions are identical. This proves equality of
the entire consequence relations, not only of their closed theorems.
\end{proof}

\subsection{An obstruction for the opposite guard orientation}
\label{compat:left-guard-obstruction}

The direction of the guard matters before zero-source transport has
been restored. The following statement explains why it is not enough
to combine the two older restrictions without changing the guard.

Consider a free-syntax calculus whose object formulas are equations and
$D$ formulas. Admit a $D$ formula with atomic matrix only when the left
side of that matrix has a fixed constructor head which survives every
capture-free numerical substitution. For example, this applies to
\[
 I(y,0;u.h)=I(y,0;u.k)
 \quad\hbox{and}\quad
 C(y,t,s)=t.
\]
Use the same computations, source-restricted transport, DE/Use,
compound case introduction, and atomic progress introduction as above.
Ordinary hypotheses are retained formulas. Clauses introduced by DE
have the corresponding admissible matrix. No arbitrary additional
clause is present in the statement below.

Call a formula \emph{safe} when it is a $D$ formula, or an equation
whose left side is not literally zero. Call a retained clause safe when
its antecedent and all its numerical instances are safe.

\begin{proposition}[Inert zero-left hypothesis]
\label{compat:inert-zero}
In this left-constructor-guard calculus, if a safe formula $F$ is
derivable from a context of safe hypotheses and safe clauses together
with one extra hypothesis $0=t$, then $F$ is derivable from the same
context with that extra hypothesis erased.
\end{proposition}
\begin{proof}
Induct on a derivation whose conclusion is safe. A safe hypothesis
cannot be the erased zero-left equation. Computation conclusions are
safe. Atomic progress introduces a local hypothesis and has a premise
with the fixed constructor head on the left; these are safe. Compound
case introduction has a compound local hypothesis and premise. In DE
the major formula is compound, and its new clause is safe because its
matrix is either compound or has the stipulated left constructor head.
Use on a safe clause has a safe antecedent, to which the induction
hypothesis applies. Thus all these cases admit erasure recursively.

For transport, its major $p=q$ is safe because $p\not\equiv0$.
If the conclusion $A(q)$ is safe, then its minor $A(p)$ is also safe.
This is immediate for a $D$ root. For an equation, a numerical template
whose left instance at the nonzero syntactic term $p$ is literally $0$
must have constant left side $0$: neither a constructor-headed term
nor an occurrence of the hole can become the free-syntax literal zero
when the hole is filled by $p$. Its instance at $q$ would then also
have left side $0$, contrary to safety. The induction hypothesis
therefore applies to both transport premises. This proves erasure.
\end{proof}

For the two concrete guarded shapes displayed above, the calculus is
sound in the standard natural numbers with the ordinary meanings of
the binders. Consequently it cannot prove $1=0$ without assumptions.
Proposition~\ref{compat:inert-zero} therefore implies
\[
                      0=1\ \nvdash\ 1=0.
\]
Thus these source-restricted presentations fail to recover even this
instance of ordinary equality symmetry or explosion. The proposition
does not apply to arbitrary externally supplied clauses with a
zero-left antecedent and a safe consequent, to definitional quotients
of term syntax, or to other additional inference rules. Those would
change its hypotheses. It applies to the proposed direct combination
of the original canonical guarded language and the source restriction.

The mirrored guard escapes the obstruction because its left term may
be the induction variable itself. The bootstrap through right expansion,
shared-left implication, and equation induction makes that extra
possibility effective.
